%! TEX root = k3_tanah_temu5.tex
% the main TeX file which is intended to compile, :VimtexReload after adjustment
% :h vimtex-tex-root
\documentclass[../main.tex]{subfiles}
% \includeonlyframes{current}
% \setbeameroption{hide notes} % Only slides
% \setbeameroption{show only notes} % Only notes
% \setbeameroption{show notes on second screen=right} % Both


\title{Keselamatan dan Kesehatan Kerja (K3) Pekerjaan Tanah \& Struktur}
\subtitle{Panduan Standar Keselamatan Konstruksi di Area Kerja (Tanah dan Struktur)}
\author{Pelatihan K3 Konstruksi}
\date{\today}

\begin{document}

% -----------------------------------------------------------------------------
% SLIDE 1: Title Slide
% -----------------------------------------------------------------------------
\begin{frame}[mytitle=center,c,mycolor=digiPH_darkblue,dark]
	\titlepage
\end{frame}

% -----------------------------------------------------------------------------
% SLIDE 2: Course Objectives & Scope
% -----------------------------------------------------------------------------
\begin{frame}[mytitle=standard,t,mycolor=digiPH_gray,light]
	\frametitle{Tujuan \& Cakupan Pelatihan}
	\framesubtitle{K3 Pekerjaan Tanah \& Struktur}
	\begin{itemize}
		\item \textbf{Tujuan utama:} Memahami potensi bahaya utama dan strategi mitigasi pada pekerjaan penggalian dan struktur.
		\item \textbf{Cakupan Topik:}
		      \begin{itemize}
			      \item K3 Pekerjaan Tanah (\textit{Excavation \& Earthworks})
			      \item K3 Pekerjaan Struktur (\textit{Formwork, Rebar, Concrete, \& Heights})
		      \end{itemize}
	\end{itemize}
\end{frame}

% -----------------------------------------------------------------------------
% SLIDE 3: Overview Pekerjaan Tanah & Karakteristik Tanah
% -----------------------------------------------------------------------------
\begin{frame}[mytitle=center,t,mycolor=digiPH_cream,light]
	\frametitle{Pekerjaan Tanah}
	\framesubtitle{Gambaran Umum \& Karakteristik Tanah}
	\begin{itemize}
		\item \textbf{Overview:} Pekerjaan galian bervariasi dari parit pipa dangkal hingga pondasi dalam; pelaksanaan yang salah berisiko cedera fatal.
		\item \textbf{Jenis Pekerjaan Tanah:} Galian (\textit{Excavation}), Timbunan (\textit{Embankment}), Pemadatan (\textit{Compaction}), dan Bawah Tanah (\textit{Underground}).
		\item \textbf{Karakteristik Tanah:} Kontrol K3 harus disesuaikan dengan jenis tanah (lempung basah/kering, cadas, pasir basah/kering, kerikil, lumpur).
	\end{itemize}
\end{frame}

% -----------------------------------------------------------------------------
% SLIDE 4: Potensi Bahaya Utama Penggalian
% -----------------------------------------------------------------------------
\begin{frame}[mytitle=standard,t,mycolor=white,light]
	\frametitle{Pekerjaan Tanah}
	\framesubtitle{Potensi Bahaya Utama Penggalian}
	\textbf{Batas Risiko:} Galian dengan kedalaman lebih dari 1,2 meter memerlukan perencanaan cermat dan penopang yang memadai.
	\vspace{0.5em}
	\begin{enumerate}
		\item Runtuhnya dinding samping (\textit{Cave-ins / Side wall collapse})
		\item Jatuhnya bahan/material mengenai pekerja (\textit{Falling materials})
		\item Terperosoknya orang dan kendaraan (\textit{Falling workers \& vehicles})
		\item Kerusakan struktur di dekatnya (\textit{Impact on nearby structures})
		\item Bahaya fasilitas bawah tanah (\textit{Underground utilities})
	\end{enumerate}
\end{frame}

% -----------------------------------------------------------------------------
% SLIDE 5: Pencegahan Runtuhnya Dinding Galian
% -----------------------------------------------------------------------------
\begin{frame}[mytitle=standard,t,mycolor=digiPH_gray,light]
	\frametitle{Pekerjaan Tanah}
	\framesubtitle{Pencegahan Runtuhnya Dinding Galian}
	\begin{itemize}
		\item \textbf{Penopang Struktural:} Cegah keruntuhan menggunakan kemiringan aman (\textit{sloping}), papan penahan (\textit{shoring}), \textit{trench box}, atau \textit{sheet piling}.
		\item \textbf{Aturan Pelaksanaan:}
		      \begin{itemize}
			      \item Pasang penopang segera seiring kemajuan penggalian.
			      \item Dilarang bekerja di atas penopang yang belum diverifikasi.
			      \item Pekerjaan wajib diarahkan oleh pengawas yang kompeten.
		      \end{itemize}
		\item \textbf{Inspeksi Wajib:} Personel terampil wajib memeriksa galian setiap hari (awal/akhir shift) dan setelah kejadian yang mempengaruhi stabilitas.
	\end{itemize}
\end{frame}

% -----------------------------------------------------------------------------
% SLIDE 6: Pencegahan Jatuh: Material, Orang, & Kendaraan
% -----------------------------------------------------------------------------
\begin{frame}[mytitle=imageplus,t,mycolor=white,light]
	\frametitle{Pekerjaan Tanah}
	\framesubtitle{Pencegahan Jatuh: Material, Orang, \& Kendaraan}
	\begin{itemize}
		\item \textbf{Bahaya Material:} Jauhkan tumpukan tanah dan peralatan minimal 0,5 m dari tepi; pasang papan perancah (\textit{toe boards}) untuk menahan puing-puing.
		\item \textbf{MANDAT HELM:} Helm keselamatan wajib digunakan untuk melindungi dari jatuhan batu/puing.
		\item \textbf{Proteksi Jatuh (Orang):} Tepi galian kedalaman $> 2\text{ m}$ memerlukan penghalang kokoh; area publik wajib dipagari penuh.
		\item \textbf{Keselamatan Kendaraan:} Gunakan \textit{stop-logs}/\textit{wheel chocks}, jaga jarak kendaraan dari tepi, dan beri cat yang jelas pada penghalang.
	\end{itemize}
\end{frame}

% -----------------------------------------------------------------------------
% SLIDE 7: Keselamatan Fasilitas Bawah Tanah & Gas Berbahaya
% -----------------------------------------------------------------------------
\begin{frame}[mytitle=standard,t,mycolor=digiPH_cream,light]
	\frametitle{Pekerjaan Tanah}
	\framesubtitle{Fasilitas Bawah Tanah \& Gas Berbahaya}
	\begin{itemize}
		\item \textbf{Deteksi Utilitas:} Periksa peta utilitas dan gunakan \textit{cable locator} sebelum menggali (catatan: alat tidak dapat mendeteksi pipa plastik).
		\item \textbf{Penggalian Aman:} Gali secara manual dengan alat tangan (sekop, bukan gancu) dalam jarak 0,5 m dari kabel/pipa.
		\item \textbf{Pengendalian Gas \& Uap:}
		      \begin{itemize}
			      \item Dilarang mengoperasikan mesin bensin/diesel di dalam parit tanpa ventilasi yang baik.
			      \item Periksa kualitas udara parit sebelum setiap shift.
			      \item Jika terjadi kebocoran pipa gas: dilarang merokok/nyala api terbuka dan segera evakuasi.
		      \end{itemize}
	\end{itemize}
\end{frame}

% -----------------------------------------------------------------------------
% SLIDE 8: Safe Access & Water Management in Excavations
% -----------------------------------------------------------------------------
\begin{frame}[mytitle=center,t,mycolor=digiPH_gray,light]
	\frametitle{Pekerjaan Tanah}
	\framesubtitle{Akses Aman \& Pengelolaan Air}
	\begin{itemize}
		\item \textbf{Akses Keluar/Masuk Aman:} Sediakan tangga yang sesuai atau jalan miring (\textit{ramp}) khusus di dalam parit galian.
		\item \textbf{Akumulasi Air:} Siapkan pompa air yang memadai untuk menjaga galian tetap kering.
		\item \textbf{Kesiapsiagaan Darurat:} Tetapkan rencana tanggap darurat dan rambu petunjuk yang jelas jika terjadi kerusakan utilitas.
	\end{itemize}
\end{frame}

% -----------------------------------------------------------------------------
% SLIDE 9: Summary Table – Mitigasi Bahaya Pekerjaan Tanah
% -----------------------------------------------------------------------------
\begin{frame}[mytitle=standard,t,mycolor=white,light]
	\frametitle{Pekerjaan Tanah}
	\framesubtitle{Rangkuman Mitigasi Bahaya}
	\centering
	\small
	\begin{tabularx}{\textwidth}{l X}
		\toprule
		\textbf{Potensi Bahaya} & \textbf{Tindakan Pencegahan Utama}                                                  \\
		\midrule
		Keruntuhan Dinding      & \textit{Sloping}, \textit{benching}, \textit{trench box}, \textit{hydraulic shores} \\
		Material Jatuh          & Jarak tanah galian $\ge 0,5\text{ m}$, \textit{toe boards}, helm K3                 \\
		Orang/Kendaraan Jatuh   & Penghalang (kedalaman $>2\text{ m}$), pemagaran, \textit{stop-blocks}               \\
		Fasilitas Bawah Tanah   & Peta utilitas, penggalian manual dalam jarak 0,5 m                                  \\
		Gas / Uap Berbahaya     & Pengujian udara, ventilasi memadai, mesin di luar galian                            \\
		\bottomrule
	\end{tabularx}
\end{frame}

% -----------------------------------------------------------------------------
% SLIDE 10: Pekerjaan Struktur: Anjungan Kerja & Perancah (Scaffolding)
% -----------------------------------------------------------------------------
\begin{frame}[mytitle=standard,t,mycolor=digiPH_gray,light]
	\frametitle{Pekerjaan Struktur}
	\framesubtitle{Anjungan Kerja \& Perancah (\textit{Scaffolding})}
	\begin{itemize}
		\item \textbf{Persyaratan Anjungan Kerja:} Lebar minimum 60 cm, bebas dari celah/lubang, bebas bahaya tersandung/tergelincir, dan dilengkapi \textit{toe boards}.
		\item \textbf{Dasar Perancah:}
		      \begin{itemize}
			      \item Dirancang, dipasang, dan dibongkar hanya oleh personel yang kompeten.
			      \item Didirikan di atas lantai keras dan rata yang mampu menahan beban (dilarang menggunakan batu bata/drum bekas).
			      \item Dirancang untuk menahan minimal $4\times$ beban maksimum yang direncanakan.
			      \item Lantai kerja perancah harus memiliki lebar minimal 46 cm.
		      \end{itemize}
	\end{itemize}
\end{frame}

% -----------------------------------------------------------------------------
% SLIDE 11: Keselamatan Perancah Menara (Tower Scaffolding)
% -----------------------------------------------------------------------------
\begin{frame}[mytitle=imageplus,t,mycolor=white,light]
	\frametitle{Pekerjaan Struktur}
	\framesubtitle{Keselamatan Perancah Menara (\textit{Tower Scaffolding})}
	\begin{itemize}
		\item \textbf{Aturan Pemasangan:} Tegak lurus/vertikal di tanah keras \& rata; kunci roda penopang saat digunakan.
		\item \textbf{Akses Internal:} Wajib menggunakan tangga internal (memanjat dari luar dapat menyebabkan perancah terbalik).
		\item \textbf{Pengaman Jatuh:} Pasang pelindung tepi (\textit{guardrails} \& \textit{toe boards}) untuk ketinggian $\ge 2\text{ m}$.
		\item \textbf{Pengikat/Jangkar Dilakukan Jika:} Menara bertingkat banyak, terpapar angin kencang, digunakan untuk penyemprotan tekanan tinggi, atau memiliki dasar yang sempit.
	\end{itemize}
\end{frame}

% -----------------------------------------------------------------------------
% SLIDE 12: K3 Pekerjaan Pembesian (Rebar Works)
% -----------------------------------------------------------------------------
\begin{frame}[mytitle=standard,t,mycolor=digiPH_cream,light]
	\frametitle{Pekerjaan Struktur}
	\framesubtitle{K3 Pekerjaan Pembesian (\textit{Rebar Works})}
	\begin{itemize}
		\item \textbf{Penanganan Besi Panjang:} Memerlukan jumlah personel yang cukup untuk mencegah pembengkokan atau risiko jatuh.
		\item \textbf{Besi Vertikal:} Amankan dengan penopang kayu/bambu sementara untuk mencegah tekukan/runtuh.
		\item \textbf{Bahaya Tertusuk:} Tutup ujung besi vertikal yang terbuka dengan \textit{protective caps} atau pelindung bambu.
		\item \textbf{Operasi Crane:} Gunakan \textit{wire rope slings} untuk mengikat besi dan pandu pengangkatan dengan isyarat peluit.
		\item \textbf{Larangan Memanjat:} Sangat dilarang memanjat naik/turun melalui besi yang terpasang; gunakan perancah.
	\end{itemize}
\end{frame}

% -----------------------------------------------------------------------------
% SLIDE 13: K3 Pekerjaan Pengecoran (Concrete Works)
% -----------------------------------------------------------------------------
\begin{frame}[mytitle=center,t,mycolor=white,light]
	\frametitle{Pekerjaan Struktur}
	\framesubtitle{K3 Pekerjaan Pengecoran (\textit{Concrete Works})}
	\begin{itemize}
		\item \textbf{Risiko Utama:} Jatuh saat pemasangan bekisting, bekisting runtuh, kejatuhan puing, debu silika, dan ketegangan otot.
		\item \textbf{Pengendalian Keselamatan:}
		      \begin{itemize}
			      \item Periksa bekisting, pengikat, dan penopang sementara sebelum pengecoran dilakukan.
			      \item Distribusikan beban beton secara merata; hindari penumpukan beban terpusat pada struktur sementara.
		      \end{itemize}
		\item \textbf{Mandat APD:} Sarung tangan \& sepatu boot karet terhadap luka bakar kimia beton basah; respirator terhadap debu silika.
	\end{itemize}
\end{frame}

% -----------------------------------------------------------------------------
% SLIDE 14: K3 Pekerjaan Shotcrete
% -----------------------------------------------------------------------------
\begin{frame}[mytitle=standard,c,mycolor=digiPH_ocean,dark]
	\frametitle{Pekerjaan Struktur}
	\framesubtitle{K3 Pekerjaan \textit{Shotcrete}}
	\begin{itemize}
		\item \textbf{Persyaratan APD untuk Operator:}
		      \begin{itemize}
			      \item Masker pelindung pernapasan (\textit{respirator})
			      \item Kacamata keselamatan pelindung debu (\textit{dust safety goggles})
			      \item Sarung tangan karet
			      \item Sepatu boot karet
		      \end{itemize}
		\item \textbf{Bahaya Kesehatan:} Kontak langsung dengan semen basah dapat menyebabkan iritasi kulit, reaksi alergi, dan luka bakar kimia.
	\end{itemize}
\end{frame}

% -----------------------------------------------------------------------------
% SLIDE 15: K3 Pekerjaan di Ketinggian (Working at Heights)
% -----------------------------------------------------------------------------
\begin{frame}[mytitle=standard,t,mycolor=digiPH_gray,light]
	\frametitle{Pekerjaan Struktur}
	\framesubtitle{K3 Pekerjaan di Ketinggian}
	\begin{itemize}
		\item \textbf{Konteks Risiko:} Jatuh menyumbang $\sim 50\%$ kematian di proyek konstruksi; perlindungan wajib untuk ketinggian $> 2\text{ m}$.
		\item \textbf{Dimensi Pelindung Tepi (\textit{Edge Protection}):}
		      \begin{itemize}
			      \item \textit{Top Guardrail}: 91 cm di atas lantai
			      \item \textit{Intermediate Rail}: Celah tidak boleh melebihi 47 cm
			      \item \textit{Toe Board}: Tinggi minimum 15 cm
		      \end{itemize}
		\item \textbf{Verifikasi Struktur:} Pastikan platform/lantai kerja mampu menahan beban pekerja dan peralatan tanpa risiko terguling.
	\end{itemize}
\end{frame}

% -----------------------------------------------------------------------------
% SLIDE 16: Checklists Keselamatan Kerja di Ketinggian & Alat
% -----------------------------------------------------------------------------
\begin{frame}[mytitle=imageplus,t,mycolor=white,light]
	\frametitle{Pekerjaan Struktur}
	\framesubtitle{Daftar Periksa Ketinggian \& Pengaman Mesin}
	\begin{itemize}
		\item \textbf{Daftar Periksa Sebelum Bekerja:}
		      \begin{itemize}
			      \item Akses masuk/keluar aman telah diverifikasi.
			      \item Semua tepi terbuka dilindungi \textit{guardrail} dan \textit{toe board}.
			      \item Peralatan diinspeksi dan pekerja telah dilatih.
			      \item \textit{Full body harness}, sarung tangan, sepatu keselamatan, dan helm digunakan dengan benar.
		      \end{itemize}
		\item \textbf{Pengaman Mesin (\textit{Machine Guarding}):} Pasang pelindung pada sabuk, puli, roda gigi, dan bagian berputar; pasang kembali segera setelah perbaikan.
	\end{itemize}
\end{frame}

% -----------------------------------------------------------------------------
% SLIDE 17: Summary & Q&A
% -----------------------------------------------------------------------------
\begin{frame}[mytitle=center,c,mycolor=digiPH_darkblue,dark]
	\frametitle{Rangkuman \& Tanya Jawab}
	\framesubtitle{K3 Pekerjaan Tanah \& Struktur}
	\begin{itemize}
		\item \textbf{Poin Utama:} Pekerjaan tanah dan struktur memiliki risiko tinggi, tetapi perencanaan sistematis, inspeksi harian, dan APD standar dapat mengeliminasi sebagian besar kecelakaan.
		\item \textbf{Tindakan Wajib:} Inspeksi area kerja setiap hari, ikuti petunjuk pengawas kompeten, dan laporkan bahaya dengan segera.
	\end{itemize}
	\vspace{1em}
	\centering
	\Large \textbf{Sesi Tanya Jawab Dibuka}
\end{frame}

\end{document}
